\section{Coverage at stretched-exponential scales}
\label{sec:coverage}

The energy and rotation estimates now supply references for the dynamical
construction. Given a target $x$ in a subset $S$ of the cube, we seek a
reference $z\in S$ with $|R(x,z)|\ge q_*$ among independent Gibbs samples.
The required estimate has two parts: a small total Gibbs mass of targets
may be excluded, and every remaining target must have a neighborhood of
sufficient Gibbs mass. Both masses are measured by the original Gibbs
law, without conditioning on $S$. Thus the neighborhood mass is the
success probability of one sample from a reference bank.

We prove the estimate with both mass bounds on the scale $L=n^\kap$.
If a set has too few overlapping pairs, fresh Gaussian disorder produces
a gain in its restricted partition sum. The energy upper tail limits
the loss from decreasing the original disorder, leaving a net increase
that the rotation estimate makes unlikely. To make this argument uniform
over all subsets, we select a violating set using the original disorder
before sampling the perturbation.

\begin{proposition}[Stretched-exponential coverage]
\label{prop:coverage}
There exists $q_*\in(0,1)$, depending only on the fixed inverse
temperature $\beta>1$, with the following property.  For each fixed
$D_1>0$, there exists a fixed $D_3=D_3(\beta,D_1)>0$ and a sequence of
events $\mathcal C_n(D_1)$, measurable with respect to the disorder,
such that $\PP_J(\mathcal C_n(D_1))\longrightarrow1$ and, on
$\mathcal C_n(D_1)$,
\begin{equation}
 \pi^J\!\left(
   \left\{x\in S:
     \pi^J\{z\in S:|R(x,z)|\ge q_*\}<e^{-D_3L}
   \right\}\right)
 <e^{-D_1L}
 \qquad\text{for every }S\subset\Sig.
 \label{eq:coverage}
\end{equation}
\end{proposition}

\begin{proof}
We first control conditional pair mass and then pass to individual targets.
Let $e_1\in(0,\beta/2)$ and $c_E>0$ be given by
Lemma~\ref{static-lem:coverage-energy}, and define
\[
 \mathcal E_n
 =\left\{\pi^J\{x:H^J(x)>ne_1\}\le e^{-c_En}\right\}.
\]
Thus $\PP_J(\mathcal E_n)\to1$.  Define
\[
 g=\frac{\beta^2}{4}-\frac{\beta e_1}{2}>0,
 \qquad
 0<q_*<1,\qquad \beta^2q_*<\frac g8.
\]
Fix this threshold once for both the present proposition and
Appendix~\ref{poly-sec:coverage}; the choice precedes $D_1$.
Given $D_1>0$, choose fixed constants
\begin{equation}
 M>\frac{8(D_1+1)}g,
 \qquad
 D_2>\frac{\beta^2M}{2}+1,
 \qquad
 t=\frac{ML}{n}.
 \label{eq:coverage-parameters}
\end{equation}
For all sufficiently large $n$, we have $0<t<1/2$;
moreover, $t\to0$ and $nt=ML\to\infty$.

\par\medskip\noindent\textbf{Pair mass in every sufficiently large set.}\enspace
Our first claim is that, with disorder probability tending to one,
\begin{equation}
 \bigl(\pi^J(\,\cdot\mid S)\bigr)^{\otimes2}
   \{|R(x,y)|\ge q_*\}
 \ge e^{-D_2L}
 \quad\text{for every }S\subset\Sig
 \text{ with }\pi^J(S)\ge e^{-D_1L}.
 \label{eq:coverage-pair}
\end{equation}
Write $\mathcal V_n$ for the failure event. Order the subsets of $\Sig$
deterministically and, for $J\in\mathcal E_n\cap\mathcal V_n$, select
the first violating member $S(J)$. The Gibbs quantities for each fixed
subset are measurable in $J$, so the selection is measurable. Once we
condition on $J$, the chosen set is fixed and cannot depend on the
independent disorder used later.

Fix such a disorder, write $\pi=\pi^J$, and remove the upper-energy
tail from the selected set:
\[
 T=S\cap\{x:H^J(x)\le ne_1\},
 \qquad \nu=\pi(\,\cdot\mid T).
\]
Since $L=o(n)$, for all sufficiently large $n$,
\[
 e^{-c_En}\le\frac12e^{-D_1L},
 \qquad
 \pi(T)\ge\frac12\pi(S)\ge\frac12e^{-D_1L}.
\]
The retained mass is at least half the original mass. Comparing the
conditional pair laws gives
\begin{equation}
 \nu^{\otimes2}\{|R(x,y)|\ge q_*\}
 \le
 \left(\frac{\pi(S)}{\pi(T)}\right)^2
 \bigl(\pi(\,\cdot\mid S)\bigr)^{\otimes2}
   \{|R(x,y)|\ge q_*\}
 <4e^{-D_2L}.
 \label{eq:coverage-trimming}
\end{equation}

\par\medskip\noindent\textbf{Gain from independent disorder.}\enspace
Sample an independent Gaussian disorder $G$ and set
\[
 W=\nu\!\left(e^{\beta\sqrt t H^G}\right),
 \qquad J'=\sqrt{1-t}\,J+\sqrt t\,G.
\]
The subscript $G$ means that the original disorder and its selected set
are held fixed. The off-diagonal normalization gives
\[
 \EE_G H^G(x)^2=\frac{n-1}{2},
 \qquad
 \EE_G[H^G(x)H^G(y)]=\frac{nR(x,y)^2-1}{2}.
\]
The Gaussian exponential moment formula consequently yields
\begin{align}
 \EE_GW&=e^{\beta^2t(n-1)/4},\notag\\
 \frac{\EE_GW^2}{(\EE_GW)^2}
 &=\nu^{\otimes2}\!\left(
      e^{\beta^2t(nR(x,y)^2-1)/2}\right)\notag\\
 &\le e^{\beta^2ntq_*^2/2}
      +4e^{(\beta^2M/2-D_2)L}
 \le 2e^{\beta^2ntq_*^2/2}
 \label{eq:coverage-moments}
\end{align}
for all sufficiently large $n$.  Split the pair sum at $|R(x,y)|=q_*$. On the large-overlap part use
\eqref{eq:coverage-trimming}, then the choice of $D_2$ in
\eqref{eq:coverage-parameters}. Paley--Zygmund consequently yields
\begin{equation}
 \PP_G\!\left(W\ge\frac12\EE_GW\right)
 \ge\frac18e^{-A_0},
 \qquad A_0=\frac{\beta^2ntq_*^2}{2}.
 \label{eq:coverage-pz}
\end{equation}

Concentration raises this lower probability to one half after a controlled
loss in the threshold. For $f(G)=\log W$, differentiation gives
\[
 \left|\frac{\partial f}{\partial G_{ij}}\right|
 \le\beta\sqrt{\frac tn},
 \qquad
 \norm{\nabla f}_2^2\le\frac{\beta^2t(n-1)}2=:\sigma^2.
\]
Thus Gaussian concentration applies to $f$ with Lipschitz constant
at most $\sigma$.  Write
\[
 a=\frac{\beta^2t(n-1)}4-\log2,
 \qquad \overline f=\EE_Gf.
\]
When $a>\overline f$, compare the upper concentration tail with
\eqref{eq:coverage-pz}:
\[
 \frac18e^{-A_0}
 \le\PP_G(f\ge a)
 \le\exp\!\left(-\frac{(a-\overline f)^2}{2\sigma^2}\right),
\]
so $\overline f\ge a-\sigma\sqrt{2(A_0+\log8)}$.
If $a\le\overline f$, the same lower bound holds immediately.
In either case the lower concentration tail yields
\[
 \PP_G\!\left(
 f\ge a-\sigma\sqrt{2(A_0+\log8)}
            -\sigma\sqrt{2\log2}
 \right)\ge\frac12.
\]
Since $\sigma\le\beta\sqrt{nt/2}$, we have
\[
 \sigma\sqrt{2(A_0+\log8)}+\sigma\sqrt{2\log2}
 \le\frac{\beta^2ntq_*}{\sqrt2}+C_\beta\sqrt{nt}
 \le\beta^2ntq_*+C_\beta\sqrt{nt}.
\]
Consequently, with conditional probability at least one half,
\begin{equation}
 \log W\ge
 \frac{\beta^2t(n-1)}4-\beta^2ntq_*
                  -C_\beta\sqrt{nt}-\log2.
 \label{eq:coverage-logw}
\end{equation}
The constants and the large-$n$ threshold are independent of the selected
disorder and its violating set.

\par\medskip\noindent\textbf{Comparison with the rotation cost.}\enspace
The gain from $G$ must be compared with the change in the original energy.
Linearity in the disorder gives
\begin{align*}
 \frac{Z^{J'}}{Z^J}
 &=\pi\!\left(
    e^{\beta(\sqrt{1-t}-1)H^J+\beta\sqrt t H^G}\right)\\
 &\ge\pi(T)e^{-\beta ne_1(1-\sqrt{1-t})}W.
\end{align*}
On $T$, the original energy is at most $ne_1$. Its negative coefficient
$\sqrt{1-t}-1$ gives the last lower bound. Combining this with
\eqref{eq:coverage-logw}, we obtain
\begin{align*}
 \log\frac{Z^{J'}}{Z^J}
 &\ge -D_1L-\log2-\beta ne_1(1-\sqrt{1-t})+\log W\\
 &\ge (g-\beta^2q_*)nt-D_1L
        -C_\beta\sqrt{nt}-C_\beta(nt^2+1).
\end{align*}
Here we used $1-\sqrt{1-t}=t/2+O(t^2)$, uniformly for
$0<t<1/2$.  Dividing the final lower bound by $nt$, its limit
inferior is at least
\[
 g-\beta^2q_* -\frac{D_1}{M}>\frac{3g}{4}.
\]
It follows that, for all sufficiently large $n$ and every
$J\in\mathcal E_n\cap\mathcal V_n$,
\begin{equation}
 \PP_G\!\left(
    \log Z^{J'}-\log Z^J\ge\frac g2nt
    \right)\ge\frac12.
 \label{eq:coverage-gain}
\end{equation}

Average the conditional gain over $J$ and use the unconditional rotation
bound \eqref{eq:rotation-tail}:
\begin{align*}
 \PP_J(\mathcal E_n\cap\mathcal V_n)
 &\le2\PP_{J,G}\!\left(
       \log Z^{J'}-\log Z^J\ge\frac g2nt\right)\\
 &\le2e^{-c'_\beta nt}
   =2e^{-c'_\beta ML}\longrightarrow0
\end{align*}
for a fixed $c'_\beta>0$.  Adding
$\PP_J(\mathcal E_n^c)=o(1)$ proves
$\PP_J(\mathcal V_n)=o(1)$, and hence the simultaneous pair
statement in Equation~\eqref{eq:coverage-pair}.  Since $\mathcal V_n$ is the event that any violating subset exists, this
proves the simultaneous claim. No union bound over subsets is needed.

\par\medskip\noindent\textbf{From pairs to individual targets.}\enspace
Choose a fixed $D_3>D_1+D_2$, and set
$\mathcal C_n(D_1)=\mathcal V_n^c$.  Fix $J$ in this event and
any $S\subset\Sig$, and let
\[
 B=\{x\in S:\pi^J\{z\in S:|R(x,z)|\ge q_*\}<e^{-D_3L}\}.
\]
Suppose that $\pi^J(B)\ge e^{-D_1L}$.  Since $B\subset S$,
\begin{align*}
 \bigl(\pi^J(\,\cdot\mid B)\bigr)^{\otimes2}
     \{|R(x,y)|\ge q_*\}
 &=\frac{1}{\pi^J(B)^2}
   \sum_{x\in B}\pi^J(x)
      \pi^J\{y\in B:|R(x,y)|\ge q_*\}\\
 &<\frac{e^{-D_3L}}{\pi^J(B)}
 \le e^{-(D_3-D_1)L}
 <e^{-D_2L}.
\end{align*}
The set $B$ itself satisfies the mass hypothesis of
\eqref{eq:coverage-pair}, so this is a contradiction. Thus
$\pi^J(B)<e^{-D_1L}$ for every $S$, proving the target-wise statement.
\end{proof}
